\documentclass[border=10pt]{standalone}
\usepackage{tikz}
\usetikzlibrary{matrix}

\begin{document}
\begin{tikzpicture}[
    header/.style={fill=green!50!black, text=white, font=\bfseries, align=left},
    data/.style={fill=green!10, align=left},
    data-right/.style={fill=green!10, align=right},
    footer/.style={fill=blue!50, text=white, font=\bfseries, align=center, text width=8cm}
]

\matrix (m) [matrix of nodes,
    nodes={minimum height=0.8cm, inner sep=4pt, outer sep=0, draw=black},
    column 1/.style={nodes=data, text width=4cm},
    column 2/.style={nodes=data-right, text width=2cm},
    column 3/.style={nodes=data-right, text width=2cm},
    row 1/.style={nodes=header},
    column sep=-\pgflinewidth,
    row sep=-\pgflinewidth,
    ]
{
    Region & Temp. Increase (°C) & Yield Variation (\%) \\
    Midwestern United States & 1.5 & $-10$ \\
    Punjab, India & 1.2 & $-5$ \\
    Western Europe & 0.9 & $-3$ \\
    Brazilian Cerrado & 2.0 & $-15$ \\
    Southeast Asia & 1.8 & $-12$ \\
    East Africa & 1.1 & $-8$ \\
};

% Footer node
\node[footer, below=0pt of m-6-1.south west, anchor=north west, minimum height=0.8cm] (footer) {Correlation Analysis: Temperature Changes and Agricultural Yield};

% Draw bottom line of the table
\draw (m-6-1.south west) rectangle (m-6-3.south east);

\end{tikzpicture}
\end{document}